\lysh{36652}{4}

\begin{alist}
\item Οι τετμημένες των κοινών σημείων της γραφικής παράστασης της συνάρτησης $g$ με τον
άξονα $x'x$ προσδιορίζονται από τις ρίζες της εξίσωσης $g(x) = 0$. Είναι:
\[g(x) = 0 \Leftrightarrow x^2 - 9 = 0 \Leftrightarrow x = -3 \text{ ή } x = 3.\]

Επομένως, τα ζητούμενα σημεία είναι τα $M(-3, 0)$ και $N(3, 0)$.

\item Είναι:
$f(3) = 4 \cdot 3 + 2 = 14 \ne 0$ και $f(-3) = 4 \cdot (-3) + 2 = -10 \ne 0$.

Επομένως, η γραφική παράσταση της $f$ δεν τέμνει τους άξονες σε κάποιο από τα σημεία $M(-3, 0)$ και $N(3, 0)$.

\item Έστω ότι υπάρχει κοινό σημείο $(x_0, 0)$ του άξονα $x'x$ στο οποίο τέμνονται οι γραφικές
παραστάσεις $C_f, C_g$ των $f, g$. Τότε ισχύει:
\[(f(x_0) = 0 \text{ και } g(x_0) = 0) \Leftrightarrow \begin{cases} 4x_0 + 2 = 0 \\ x_0^2 - 9 = 0 \end{cases} \Leftrightarrow \begin{cases} x_0 = -\dfrac{1}{2} \\ x_0 = \pm 3 \end{cases}\]
που είναι άτοπο. Άρα οι $C_f, C_g$ δεν έχουν κοινό σημείο πάνω στον άξονα $x'x$.

Έστω ότι οι $C_f, C_g$ έχουν κοινό σημείο πάνω στον άξονα $y'y$. Τότε έχουμε:
\[f(0) = g(0) \Leftrightarrow 2 = -9\]
που είναι άτοπο. Άρα οι γραφικές παραστάσεις συναρτήσεων $f, g$ δεν έχουν κοινό σημείο πάνω στον άξονα $y'y$.

\item Η γραφική παράσταση της $h$ είναι ευθεία, οπότε ο τύπος της είναι της μορφής
$h(x) = \alpha x + \beta, \alpha, \beta \in \mathbb{R}$. Η ευθεία διέρχεται από το σημείο $A(0, 3)$, οπότε έχουμε:
\[h(0) = 3 \Leftrightarrow \alpha \cdot 0 + \beta = 3 \Leftrightarrow \beta = 3.\]

Επομένως έχουμε $h(x) = \alpha x + 3, \alpha \in \mathbb{R}$. Επιπλέον, η γραφική παράσταση της $h$ τέμνει την γραφική παράσταση της $g$ σε σημείο του ημιάξονα $Ox$ οπότε το σημείο είναι το $N(3,0)$, άρα ισχύει:
\[h(3) = g(3) = 0 \Leftrightarrow 3\alpha + 3 = 0 \Leftrightarrow \alpha = -1.\]
Τελικά η συνάρτηση $h$ έχει τύπο $h(x) = -x + 3$.
\end{alist}

